\section{Newman-time regularity through collisions}\label{sec:newman-time}

Gaussian observation time and Newman time play different roles.
The former smooths a fixed divisor; the latter moves the divisor by
$\partial_\lambda H=-H_{zz}$. We retain the compact physical window
and fixed observation regime \eqref{eq:fixed-regime} throughout
$0\le\lambda\le1/2$. The result in this section is the full marked-pair
time modulus of \cite{OCTime,OCTimeReview}; it supplies conditional
transport of initial data, without assuming a positive minimum root gap.

\subsection{Complete-force inputs and cancellation of the pair singularity}

We use the canonical strip $|y_j|\le1$ and the following coarse input:
there is a finite absolute $A\ge1$ such that, throughout this interval,
\begin{equation}\label{eq:time-input}
 \left|\frac{H_\lambda'}{H_\lambda}(u+4\ii)\right|
      \le A\log(2+|u|),\qquad
 \#\{j:u\le x_j\le u+1\}\le A\log(2+|u|).
\end{equation}
The input is obtained from the Polymath asymptotic estimates and the
endpoint argument recorded in \cite{PolymathHeat2019,ParentModulus}.
At time zero the safe-line bound follows by locally uniform convergence
at each fixed $u$; the constant is uniform in $u$ before the limit.
For counts, enclose a unit interval by a slightly enlarged zero-free
rectangle, use the argument principle for small positive times, and
absorb the fixed enlargement into $A$. Reflection covers negative
heights and compactness covers the bounded-height remainder. No
numerical value of $A$ is asserted here.

For $X\ge10$ let $L_X=\log(2+X)$,
$\mathcal C_X=\{j:|x_j-X|\le3\}$, and let $m_X$ count labels in
$|x_j-X|\le1$. The input gives
\[
 n_X:=|\mathcal C_X|\le12AL_X,\qquad m_X\le2AL_X.
\]
Write $F_\lambda=H_\lambda'/H_\lambda$ and $w=x+4\ii$.
The complete exterior force at an active point $z=x+\ii y$ is
\begin{equation}\label{eq:exterior-force}
 E_X(z)=F_\lambda(w)-\sum_{k\in\mathcal C_X}\frac1{w-z_k}
  +\sum_{k\notin\mathcal C_X}
       \left(\frac1{z-z_k}-\frac1{w-z_k}\right).
\end{equation}
This retains the opposite-paired canonical logarithmic derivative.
The last series is absolutely convergent: its summands are at most
$5/|x-x_k|^2$, while the counts grow logarithmically. Unit shells
starting at distance two bound this series by $25AL_X$; the first
two terms cost at most $2AL_X$ and $n_X/3\le4AL_X$. Thus
\begin{equation}\label{eq:exterior-bound}
 |E_X(z)|\le31AL_X\qquad(|x-X|\le1,\ |y|\le1).
\end{equation}
For example the shell bound uses
$10A\sum_{h\ge2}[L_X+\log(h+2)]h^{-2}$;
$\sum_{h\ge2}h^{-2}\le1$ and the decreasing-function integral bound
\[
 \sum_{h\ge2}\frac{\log(h+2)}{h^2}
 \le\int_1^\infty\frac{\log(t+2)}{t^2}\dd t
 =\tfrac32\log3<3<\tfrac32L_X
\]
give the displayed constant. These estimates are only upper
allowances, not signed exterior information.

Let $f$ be a real $C^2$ test supported in $|x-X|\le1$, and put
$F_1=\sup_{|y|\le1}|\nabla f|$,
$F_2=\sup_{|y|\le1}\|D^2f\|_{\rm op}$. At a simple active zero,
implicit differentiation gives
\[
 \dot z_j=2\sum_{k\in\mathcal C_X,\,k\ne j}\frac1{z_j-z_k}
                  +2E_X(z_j).
\]
For $d=x_j-x_k$, $e=y_j-y_k$, the two ordered force terms combine
into $-2K_f(z_j,z_k)$, where
\begin{align}
 K_f(z_j,z_k)
  &=\frac{e(f_y(z_j)-f_y(z_k))-d(f_x(z_j)-f_x(z_k))}{d^2+e^2}\nonumber\\
  &=\int_0^1\frac{e^2 f_{yy}-d^2 f_{xx}}{d^2+e^2}
                     (z_k+t(z_j-z_k))\dd t.\label{eq:chord-cancellation}
\end{align}
The mixed derivatives cancel, and the chord stays in the strip.
Hence $|K_f|\le F_2$ even as the separation tends to zero.
Only pairs with at least one active label contribute, at most
$m_Xn_X$ unordered pairs. The exterior costs $2m_XF_1\,31AL_X$.
We have proved, at simple times,
\begin{equation}\label{eq:compact-time-test}
 \left|\partial_\lambda\sum_j f(z_j)\right|
       \le A^2L_X^2(48F_2+124F_1).
\end{equation}
This is the useful cancellation: bounding the two velocities separately
would introduce a reciprocal root separation that the paired test does
not require.

\subsection{The full pair and connected moduli}

Set $f_{X,u}(z)=JW_X(z)G_\epsilon(X+\ell u-z)$ and define
$r_{\lambda,X},d_{\lambda,X},R_{2,\lambda,X}$ as before.
On the canonical strip the transverse cutoff is identically one,
including its first two derivatives at the boundary. The horizontal
window has derivative bounds $15/4$ and $24$. The Gaussian has
gradient bound $A_\epsilon/(\epsilon\sqrt{\ee})$ and Hessian operator
bound $A_\epsilon/\epsilon^2$. Define
\begin{align}
 a_0&=\frac{J}{2\pi\epsilon^2}=\frac8\pi,&
 a_1&=a_0\left(\frac{15}4+\frac1{\epsilon\sqrt{\ee}}\right),\nonumber\\
 a_2&=a_0\left(24+\frac{15}{2\epsilon\sqrt{\ee}}+\frac1{\epsilon^2}\right),&
 L&=\log(2+2T),\quad R_*=2ALa_0,\nonumber\\
 L_1&=A^2L^2(48a_2+124a_1),&
 L_d&=A^2L^2[96(a_0a_2+a_1^2)+248a_0a_1],\nonumber\\
 L_2&=2R_*L_1+L_d,&L_C&=4R_*L_1+L_d.
                                                        \label{eq:time-constants}
\end{align}

\begin{theorem}[Full marked-pair time modulus]\label{thm:time-modulus}
For $\lambda,\mu\in[0,1/2]$, uniformly in $X\in[T,2T]$,
\[
 \begin{aligned}
 \|r_{\lambda,X}-r_{\mu,X}\|_\infty&\le L_1|\lambda-\mu|,\\
 \|d_{\lambda,X}-d_{\mu,X}\|_\infty&\le L_d|\lambda-\mu|,\\
 \|R_{2,\lambda,X}-R_{2,\mu,X}\|_\infty&\le L_2|\lambda-\mu|.
 \end{aligned}
\]
The same first and pair bounds hold after center averaging, and
$\|C_\lambda-C_\mu\|_\infty\le L_C|\lambda-\mu|$.
The norms are on $Q$ or $Q^2$, and the estimates include multiple-zero times.
\end{theorem}
\begin{proof}
The product rule gives $|f_{X,u}|\le a_0$,
$|\nabla f_{X,u}|\le a_1$, $\|D^2f_{X,u}\|\le a_2$.
The self test $f_{X,u}f_{X,v}$ has gradient at most $2a_0a_1$ and
Hessian at most $2a_0a_2+2a_1^2$. Formula
\eqref{eq:compact-time-test} proves $L_1,L_d$ at simple times.
Also $r_{\lambda,X}\le R_*$. The raw product costs $2R_*L_1$;
subtracting the self term gives $L_2$. The connected statistic subtracts
a second product, now of averaged means, giving $L_C$.

It remains to justify integration through events. The complete family
is analytic in $(\lambda,z)$ and is not identically zero on a slice,
since $H_\lambda(0)>0$. In the padded compact rectangle
$[T-2,2T+2]+\ii[-2,2]$, local zero-free boundaries and Rouch\'e's theorem
give continuous finite multisets of roots, with multiplicity preserved.
The weighted first, self and pair sums are therefore jointly continuous.

At a zero of multiplicity $q$ at $(\lambda_0,z_0)$, put
$t=\lambda-\lambda_0$, $\zeta=z-z_0$. The heat equation makes the
lowest parabolic-degree Taylor polynomial a nonzero multiple of
\[
 \ee^{-t\partial_\zeta^2}\zeta^q
  =(2t)^{q/2}\operatorname{He}_q(\zeta/\sqrt{2t}).
\]
The polynomial identity follows by expanding the finite differential
series, or the Hermite generating function \cite{DLMFHermiteOC}.
The $q$ Hermite zeros are real and distinct: the Rodrigues expression
$\operatorname{He}_q(x)=(-1)^q\ee^{x^2/2}(\dd/\dd x)^q\ee^{-x^2/2}$
and integration by parts give orthogonality to every polynomial of
degree below $q$ for the positive Gaussian weight. If there were fewer
than $q$ real sign changes, their linear factors would form such a
polynomial whose product with $\operatorname{He}_q$ has constant sign,
contradicting orthogonality. Rescaling and applying Rouch\'e
near each one shows that all nearby roots are simple for sufficiently
small nonzero $t$; their displacements are
$\sqrt{2t}\,h_j+O(|t|)$. Thus multiple events are isolated in this
compact space-time set and there are finitely many. Integrate the
uniform simple-time estimates on the intervening intervals and use
continuity at their endpoints. The interval lengths sum to
$|\lambda-\mu|$, so no event-count factor appears. Averaging over
the unchanged center law preserves the estimates.
\end{proof}

\subsection{Exterior truncation and conditional transport}

Truncate only the absolutely convergent difference series in
\eqref{eq:exterior-force}, retaining terms with $|x-x_k|<M$, for
integer $M\ge2$. Unit counts and a decreasing-function integral give
\begin{equation}\label{eq:exterior-tail}
 e_M(X)=\frac{10A}{M-1}[L_X+\log(M+1)+1].
\end{equation}
Indeed the shell majorant is
$10A[L_X+\log(h+2)]/h^2$; its integral from $M-1$ to infinity
is bounded by \eqref{eq:exterior-tail}, using
$\log(1+s)\le s$ after rescaling. The induced scalar derivative
error is at most $2m_XF_1e_M(X)$.
Let $m_*=2AL$, and replace $L_X$ by $L$ in $e_M$ to define $e_M^*$.
The full-pair and connected derivative errors are at most, respectively,
\begin{align}
 B_{2,M}&=(4R_*m_*a_1+4m_*a_0a_1)e_M^*,\nonumber\\
 B_{C,M}&=(8R_*m_*a_1+4m_*a_0a_1)e_M^*.
                                                        \label{eq:pair-exterior-error}
\end{align}
The first coefficient pays two first-intensity derivative errors;
the second pays the self-product derivative. Connected subtraction
doubles the first coefficient. Uniform convergence allows time
integration, giving errors $B_{2,M}|\lambda-\mu|$ and
$B_{C,M}|\lambda-\mu|$. Root locations and undifferentiated intensities
remain those of the exact divisor; this is not an error theorem for
an independently truncated root dynamics.

If a function $F_0$ approximates $\overline R_{2,\mu}$ within $e_0$,
then
\begin{equation}\label{eq:initial-transport}
 \|\overline R_{2,\lambda}-F_0\|_\infty
       \le e_0+L_2|\lambda-\mu|.
\end{equation}
The connected version uses $L_C$. For large $T$, the displayed allowances
are $L_1,L_d=O_A((\log T)^4)$ and $L_2,L_C=O_A((\log T)^5)$.
They supply no initial profile and no favorable sign.

Lipschitz continuity is compatible with a kink at collision. For
$p_\lambda(z)=z^2-2\lambda$ and constant local window, the ordered
two-label density at observations $u=v=0$ is exactly
\[
 2(2\pi\epsilon^2)^{-2}\ee^{-2|\lambda|/\epsilon^2}.
\]
Its right-minus-left slope is
$-8(2\pi\epsilon^2)^{-2}/\epsilon^2$. At $\lambda=0$ both labels
remain and the factorial coefficient is $2$, not zero. This local
heat control proves neither a canonical collision nor differentiability
through one; it shows why the theorem claims a modulus rather than a
globally smooth closed equation.
